\section{1830 Gauss oraz Farkas i János Bolyai} % lang-pl
\label{sekcja_gauss_bolyai}

Carl Friedrich Gauss (1777--1855) zapozna się w Getyndze, około 1796--1799 roku, z Farkasem Bolyaiem, i ci dwaj zostaną bliskimi przyjaciółmi. % lang-pl
Farkas Bolyai (1775--1856), urodzony w Bólyi, a zmarły w Marosvásárhely, przez całe życie będzie zajmował się podstawami geometrii i aksjomatem równoległych. % lang-pl
\index[persons]{Gauss, Carl Friedrich}%
\index[persons]{Bolyai, Farkas}%

Jego syn János (1802--1860), urodzony w Kolozsvárze, mimo ostrzeżeń ojca weźmie się za ten sam problem. % lang-pl
Farkas napisze mu, że przeszedł przez tę „bezdenną noc'', która zgasiła w jego życiu wszelkie światło i radość. % lang-pl
János nie usłucha, a w 1823 roku napisze do ojca z ekscytacją, że odkrył rzeczy tak cudowne, iż sam się zdumiał, i że „z niczego stworzył dziwny nowy świat''. % lang-pl
\index[persons]{Bolyai, János}%

W 1832 roku wyniki Jánosa ukażą się jako dodatek do pracy ojca \emph{Tentamen}: Farkas, początkowo niechętny, dojdzie do entuzjazmu i sam przekona syna do publikacji. % lang-pl
Gauss odpowie z uznaniem, nazywając Jánosa geniuszem pierwszej rangi, ale doda, że sam doszedł do podobnych wyników wiele lat wcześniej -- co, jak się mówi, wywoła napięcie między ojcem i synem. % lang-pl
Później János dowie się, że już w latach 1829--1830 podobne wyniki opublikował Łobaczewski (zobacz sekcję \ref{sekcja_lobaczewski}). % lang-pl

W 1829 roku János Bolyai, Carl Friedrich Gauss oraz (niezależnie od nich) Nikołaj Łobaczewski stworzą geometrię hiperboliczną, czyli pierwszą w pełni rozwiniętą geometrię nieeuklidesową, w której piąty postulat Euklidesa nie obowiązuje. % lang-pl
Ściślej: Łobaczewski opublikuje swoje wyniki w latach 1829--1830, Bolyai w 1832 roku, a Gauss -- choć według Evesa pierwszy dojdzie do głębokich wniosków -- nie opublikuje niczego, więc honor odkrycia trzeba będzie z nimi podzielić \cite[s. 287]{eves1_1972}. % lang-pl

Gauss, według Evesa, podejdzie do problemu tak samo jak pozostali dwaj: przez trzy możliwości dla aksjomatu Playfaira -- przez punkt przechodzi więcej niż jedna, dokładnie jedna albo żadna prosta równoległa do danej -- a trzecią łatwo wykluczy przy założeniu nieskończoności prostej \cite[s. 287]{eves1_1972}. % lang-pl
On pierwszy nada nazwę \emph{geometria nieeuklidesowa}, początkowo mając na myśli własną geometrię hiperboliczną; w 1824 roku napisze też list do Taurinusa (zobacz sekcję \ref{sekcja_schweikart_taurinus}). % lang-pl

János Bolyai zbuduje geometrię, która obejmuje jednocześnie euklidesową i hiperboliczną, zależnie od parametru $k$; nazwie ją \emph{geometrią absolutną}. % lang-pl
Zauważy też, że samym rozumowaniem matematycznym nie da się rozstrzygnąć, czy geometria świata fizycznego jest euklidesowa, czy nie (zobacz sekcję \ref{sekcja_fizyka}). % lang-pl
Eves nazywa geometrią absolutną wszystko, co da się udowodnić bez piątego postulatu -- czyli naszą geometrię neutralną; do niej należy pierwsze dwadzieścia osiem twierdzeń pierwszej księgi Euklidesa \cite[s. 288]{eves1_1972}. % lang-pl
\index{geometria!absolutna} % lang-pl

Dodatek Jánosa zakończy się wynikiem o wielokątach, który Bolyai przywoła z podziwem jako „teorię wielokątów sławnego Gaussa''; Gauss nie odwzajemni się i nie pomoże Bolyaiowi w karierze, bo obawiać się będzie „wrzasku Beotów'' (czyli reakcji tych, którzy nie zrozumieją geometrii nieeuklidesowej), a po latach Bolyai ostro skrytykuje go za tę bojaźliwość \cite[s. 213]{greenberg_2010}. % lang-pl
Bolyai poda też konstrukcję cyrklem i linijką równoległej granicznej w płaszczyźnie hiperbolicznej \cite[s. 210]{greenberg_2010}. % lang-pl

% Uwaga: nie opisuję listów Gaussa (do Taurinusa 1824, do Farkasa 1832) dokładnie -- potrzebne źródło z datami.
% Źródła: Eves s. 287; Wikipedia: Nikolai Lobachevsky, Carl Friedrich Gauss.
% https://en.wikipedia.org/wiki/Nikolai_Lobachevsky
% Uwaga: Wikipedia podaje rok urodzenia 1792; Eves (s. 287) i Hartshorne (s. 305) -- 1793.

